% ===================================================================
%  TABELL Nr 14 — Gatan. — Handlarna.
%  Traduction 2026 d'apres texto/io, controlee sur texto/fr et
%  texto/nl. Consigne : voir texto/sv/10-tabell-01.tex.
%
%  LE TABLEAU LE PLUS CHARGE DU LIVRET : quatre-vingt-dix-neuf
%  renvois, et le suedois y rencontre son genitif prepose deux fois.
%  « la aquo-spricilo (49) di la fonteno (48) » et « l'emperialo (59)
%  dil omnibuso (58) » deviendraient, en genitif, « fontanens
%  vattenstrale » et « omnibussens imperial » — l'ordre s'inverse.
%  On garde donc la preposition : « vattenstralen (49) fran fontanen
%  (48) », « imperialen (59) pa omnibussen (58) ». Le suedois
%  l'admet sans peine ; c'est le meme remede qu'au tableau 7.
% ===================================================================

%%K t14-apar-1 apar
\VUsaut{4.0mm}
\VUtitre{15.0pt}{60}{TABELL Nr 14}
\VUsaut{3.0mm}

%%K t14-tit-1 sub
\VUpk{11.4pt}{40}{Gatan. --- Handlarna.}
\VUsaut{3.0mm}

%%K t14-01-1 p
1. --- Min vän Johannes, som bor på landet, kom och tillbringade tre dagar hos min familj.
\VUgras{Tills} dess hade han inte haft tillfälle att se en stor stad; därför använder vi alla
våra dagar till att promenera, och jag förklarar för honom det som han inte känner till.

%%K t14-02-1 p
2. --- Först gick vi och besökte \VUgras{observatoriet} \textsuperscript{(1)}, som
intresserade honom mycket. När vi kom tillbaka genom \VUgras{gatan} \textsuperscript{(3)} där
det \VUgras{turkiska badet} \textsuperscript{(2)} ligger, mötte vi en \VUgras{begravning}
\textsuperscript{(4)}. Framför
\VUgras{likföljet} gick \VUgras{prästerskapet}, som föregick
\VUgras{likvagnen}, i vilken \VUgras{kistan} var ställd. Många vänner följde den
\VUgras{sörjande} familjen.

%%K t14-02-2 p
Sedan följet hade gått förbi, gick vi över gatan framför det stora
\VUgras{ölhuset} \textsuperscript{(5)} i vilket många \VUgras{gäster} drack
\VUgras{öl} på terrassen, skyddade av den glasade \VUgras{markisen}
\textsuperscript{(6)}.

%%K t14-03-1 p
3. --- Johannes blev mycket förvånad över den \VUgras{sköld} som var satt mellan
\VUgras{likörhandlarens} \textsuperscript{(7)} bod och
\VUgras{musikhandlarens} \textsuperscript{(10)}. Han frågade mig om dess
betydelse; jag svarade att den utmärker det engelska \VUgras{konsulatet}
\textsuperscript{(8)}, och gjorde honom uppmärksam på \VUgras{konsuln} \textsuperscript{(9)}
som står på balkongen.

%%K t14-tit-2 sub
\VUpk{10.2pt}{40}{Utforskning av butikerna.}
\VUsaut{3.0mm}

%%K t14-04-1 p
4. --- Naturligtvis stannade vi framför utställningen av
\VUgras{musikinstrumenten}, \VUgras{harmonierna}, \VUgras{orglarna} och
\VUgras{pianona}; vi betraktade de illustrerade omslagen till
\VUgras{partituren} och till \VUgras{musikstyckena} för piano, och beundrade
\VUgras{lyran} av förgyllt trä, som används som skylt.

%%K t14-05-1 p
5. --- Dammolnen som \VUgras{gatsoparna} \textsuperscript{(11)} och
\VUgras{sopvagnen} \textsuperscript{(12)} åstadkom tvingade oss att snabbt gå
förbi \VUgras{möbelhandlarens} \textsuperscript{(13)} vackra
\VUgras{möblemang} och förbi utställningen av \VUgras{kaminer} och
\VUgras{lampor} hos \VUgras{bleckslagaren} \textsuperscript{(14)}.

%%K t14-06-1 p
6. --- För övrigt blev vi på det stället tvungna att lämna trottoaren, för den var belamrad
med packor, som man tog ur en \VUgras{flyttvagn} \textsuperscript{(16)} för att ställa dem i
ett \VUgras{magasin} \textsuperscript{(15)} som man just hade hyrt ut.

%%K t14-06-2 p
Det hindrade oss från att se \VUgras{vagnmakarens} \textsuperscript{(17)} vackra vagnar, men
hindrade oss inte från att stanna för att se på
\VUgras{packaren} \textsuperscript{(19)} som gjorde en låda åt sin granne,
\VUgras{handlaren i cyklar och bilar} \textsuperscript{(18)}. Sedan, eftersom
det redan var något sent, vände vi hem igen.

%%K t14-tit-3 sub
\VUpk{10.2pt}{40}{Promenad genom staden.}
\VUsaut{3.0mm}

%%K t14-07-1 p
7. --- Följande dag föreslog min mor Johannes att man skulle fotografera honom, sedan gav hon
mig i uppdrag att följa honom till \VUgras{fotografen} \textsuperscript{(20)}. Efter frukosten
gick vi in i konstnärens
\VUgras{ateljé}, där vi blev tvungna att vänta ganska länge. För att sysselsätta
oss betraktade vi \VUgras{porträtten} \textsuperscript{(22)} som hänger på väggen.

%%K t14-07-2 p
När det var vår tur, varade arbetet inte länge, och genast när min vän hade slutat posera
framför \VUgras{fotografiapparaten} \textsuperscript{(21)}, gick vi ner.

%%K t14-08-1 p
8. --- På första våningen såg vi genom \VUgras{frisörens} \textsuperscript{(23)} dörr, som
står öppen, hans gesäll som klippte håret på en kund.

%%K t14-08-2 p
När vi hade kommit ut på gatan, gick vi förbi \VUgras{tobaksboden} \textsuperscript{(24)}.
Johannes blev så road av den röda \VUgras{lyktan} \textsuperscript{(25)} och den
\VUgras{tjocka pipan} som används som
\VUgras{skylt} \textsuperscript{(26)}, att han stötte ihop med två gamla herrar
som samtalade medan de \VUgras{tog snus} \textsuperscript{(27)}.

%%K t14-tit-4 sub
\VUpk{10.2pt}{40}{Konditoriet.}
\VUsaut{3.0mm}

%%K t14-09-1 p
9. --- \VUgras{Sedlarna} och \VUgras{mynten} från alla länder som var utställda i
\VUgras{växlarens} \textsuperscript{(29)} \VUgras{skyltfönster} drog några ögonblick till sig
vår uppmärksamhet, men eftersom det var dags att äta ett mellanmål, gick vi in hos
\VUgras{konditorn} \textsuperscript{(31)} utan att stanna längre.

%%K t14-10-1 p
10. --- När vi såg alla de \VUgras{läckerheter} som boden innehöll,
\VUgras{tårtor, honungskakor, skorpor} och alla slags \VUgras{karameller},
kunde vi inte lätt välja. Till slut valde Johannes \VUgras{gräddtårtorna}, och jag tog bara en
liten \VUgras{körsbärstårta}, för det söta \VUgras{gör ont i mina tänder}, och jag önskar alls
inte gå och besöka \VUgras{tandläkaren} \textsuperscript{(30)} alltför ofta.

%%K t14-tit-5 sub
\VUpk{10.2pt}{40}{Maskinskrivning.}
\VUsaut{3.0mm}

%%K t14-11-1 p
11. --- För att visa min vän en \VUgras{skrivmaskin} som han hade hört talas om, men som han
inte kände till, gick vi in på \VUgras{kontoret} \textsuperscript{(32)} hos min \VUgras{kusin}
\textsuperscript{(34)}. Vi fann honom dikterande sina brev för sin \VUgras{sekreterare}
\textsuperscript{(35)} som är \VUgras{stenograf}. Knappt var ett brev färdigt, förrän en
\VUgras{maskinskrivare} \textsuperscript{(33)} skrev av det med sin maskin.
Eftersom vi inte ville avbryta min släktings sysslor, stannade vi bara några ögonblick hos
honom.

%%K t14-12-1 p
12. --- Under min kusins kontor bor en stor \VUgras{urmakare-juvelerare}
\textsuperscript{(37)} som dessutom är guldsmed. Hans praktfulla bod innehåller många
\VUgras{ur, pendyler, fickur}, \VUgras{kedjor} och dyrbara
\VUgras{smycken}, till exempel \VUgras{pärlhalsband}, \VUgras{armband} av guld,
\VUgras{örhängen} och \VUgras{fingerringar} prydda med \VUgras{ädelstenar} och
\VUgras{diamanter}. Ett av skyltfönstren är särskilt reserverat för
\VUgras{guldarbetena} och innehåller en betydande samling \VUgras{matbestick}
av silver.

%%K t14-13-1 p
13. --- När vi vek om gathörnet, märkte jag att man hade bytt ut den
\VUgras{skylt} \textsuperscript{(36)} som är satt bredvid husets
\VUgras{nummer}. Jag var i färd med att säga min anmärkning högt, när de
glada tonerna av \VUgras{den} militära musiken hördes. Vi började springa regementet till
mötes, utan att tänka på att \VUgras{det} skulle passera framför
\VUgras{hattmakarens} \textsuperscript{(39)} och \VUgras{handskmakarens}
\textsuperscript{(40)} bodar, och att vi skulle ha varit lika väl placerade för att se dess
defilering om vi hade stannat framför \VUgras{optikerns} \textsuperscript{(38)} skyltfönster,
som är helt försett med
\VUgras{pincenéer, lornjetter} och \VUgras{kikare}.

%%K t14-tit-6 sub
\VUpk{10.2pt}{40}{Marschdefilering.}
\VUsaut{3.0mm}

%%K t14-14-1 p
14. --- Framför \VUgras{regementet} \textsuperscript{(41)} gick
\VUgras{trumslagarmajoren} \textsuperscript{(42)} följd av \VUgras{trumslagarna}
\textsuperscript{(43)}, \VUgras{hornblåsarna} \textsuperscript{(44)} och hela
\VUgras{musikkåren}. \VUgras{Generalen} \textsuperscript{(45)} som var på torget
med sin adjutant, hade stannat nära \VUgras{vattenstrålen} \textsuperscript{(49)} från
\VUgras{fontänen} \textsuperscript{(48)}, för att bevista \VUgras{defileringen}.

%%K t14-14-2 p
\VUgras{Stabsofficerarna} \textsuperscript{(46)}, \VUgras{översten} och
\VUgras{överstelöjtnanten} hälsade honom med värjan. \VUgras{Majorerna} stod
framför sina \VUgras{bataljoner} \textsuperscript{(47)} och varje
\VUgras{kapten} förde befäl över sitt \VUgras{kompani}, medan
\VUgras{löjtnanterna}, \VUgras{underlöjtnanterna} och \VUgras{underofficerarna}
intog sin vanliga plats bredvid sina män.

%%K t14-15-1 p
15. --- Många förbipasserande hade samlats i \VUgras{gatukorsningen} \textsuperscript{(50)};
handlarna, \VUgras{paraplymakaren} \textsuperscript{(51)}, \VUgras{färgaren}
\textsuperscript{(53)},
\VUgras{vapensmeden} \textsuperscript{(54)} kom ut för att se
\VUgras{soldaterna}. Alla fordon, \VUgras{hyrvagnarna} \textsuperscript{(52)},
\VUgras{herrekipagen} \textsuperscript{(57)} och även \VUgras{sprutvagnen}
\textsuperscript{(56)}, ställde sig längs trottoaren.

%%K t14-tit-7 sub
\VUpk{10.2pt}{40}{Scener på gatan.}
\VUsaut{3.0mm}

%%K t14-15-2 p
En \VUgras{handelsresande} \textsuperscript{(55)} som bar sina
\VUgras{provkoffertar} i handen, gick snabbt över gatan, och de personer som satt
på \VUgras{imperialen} \textsuperscript{(59)} på \VUgras{omnibussen} \textsuperscript{(58)}
vände sig om för att njuta av skådespelet. En stackars
\VUgras{gatusångare} \textsuperscript{(60)} med sitt \VUgras{positiv}
\textsuperscript{(61)} måste avbryta sin \VUgras{visa}, för larmet av trummorna överröstade
hans stämma.

%%K t14-16-1 p
16. --- När regementet kom fram till \VUgras{tyghandeln} \textsuperscript{(64)} lämnade
\VUgras{sömmerskorna} \textsuperscript{(63)} och
\VUgras{skräddarna} \textsuperscript{(62)} sina \VUgras{symaskiner} och sitt
arbete för att se ut genom fönstret. \VUgras{Handlaren} \textsuperscript{(65)}, som just hade
satt nya \VUgras{kostymer} \textsuperscript{(66)} i sin \VUgras{skyltning}, måste låta bära in
\VUgras{tygerna} \textsuperscript{(67)} och \VUgras{linnevarorna} som var
uppställda på trottoaren, nära \VUgras{kloaköppningen} \textsuperscript{(68)}, av fruktan att
folkmassan skulle välta dem, till och med en gammal
\VUgras{gårdfarihandlare} \textsuperscript{(69)}, av vilken en portvakterska
prutade på strumpor, blev tvungen att gå in i en korridor för att avsluta sin försäljning.

%%K t14-16-2 p
Vi följde regementet ända till kasernen, \VUgras{och}, mycket nöjda med vår dag, kom vi
tillbaka vid middagens timme.

%%K t14-tit-8 sub
\VUpk{10.2pt}{40}{Olika näringar och yrken.}
\VUsaut{3.0mm}

%%K t14-17-1 p
17. --- Följande dag föreslog min far oss att besöka den lokala tidningens tryckeri. Vi antog
det med entusiasm.

%%K t14-17-2 p
\VUgras{Tryckaren} är också \VUgras{bokhandlare} \textsuperscript{(70)}
\VUgras{(= bokförsäljare)}, \VUgras{förläggare}, \VUgras{pappershandlare} och
\VUgras{bokbindare}. Han har på första våningen inrättat tidningens
\VUgras{redaktion} \textsuperscript{(71)}, och även \VUgras{gravyrateljén}, i
vilken \VUgras{litograferna} \textsuperscript{(72)} och
\VUgras{kopparstickarna} \textsuperscript{(73)} arbetar, slutligen
\VUgras{sätteriet}, i vilket \VUgras{tryckerigesällerna} \textsuperscript{(74)}
är. Det egentliga \VUgras{tryckeriet} \textsuperscript{(75)},
\VUgras{tryckpressarna} och de olika maskinerna är på bottenvåningen.

%%K t14-tit-9 sub
\VUpk{10.2pt}{40}{Johannes ställer mycket många frågor.}
\VUsaut{3.0mm}

%%K t14-18-1 p
18. --- När vi gick ut ur tryckeriet, stannade Johannes mycket förvånad framför en liten
glasad låda som var satt på en pelare, mitt emot
\VUgras{sybehörsaffären} \textsuperscript{(77)} och frågade vad den används
till. Min far förklarade för honom, att det är en \VUgras{brandvarnare}
\textsuperscript{(76)}. För övrigt var allting, i staden, för min vän ett under, från den
enkla \VUgras{stenfontänen} \textsuperscript{(78)} och den lätta \VUgras{budtrehjulingen}
\textsuperscript{(80)} som körs av en
\VUgras{piccolo} \textsuperscript{(79)} i livré, till \VUgras{postvagnen}
\textsuperscript{(81)}.

%%K t14-19-1 p
19. --- Medan min far några ögonblick lämnade oss för att tala med
\VUgras{skatteuppbördsmannen} \textsuperscript{(83)}, som hade stannat för att
samtala med en \VUgras{advokat} \textsuperscript{(82)} och en
\VUgras{notarie} \textsuperscript{(84)}, roade vi oss med att följa trafiken på
gatan med ögonen. De mest olikartade personer gick förbi oss : en
\VUgras{konditorgesäll} \textsuperscript{(85)}, som i en korg bar en praktfull
\VUgras{pastej}, kom bakom en \VUgras{klädhandlare} \textsuperscript{(86)}; en
\VUgras{lykttändare} \textsuperscript{(87)} samtalade med en
\VUgras{gasarbetare} \textsuperscript{(88)}, en \VUgras{nunna}
\textsuperscript{(89)} som höll en föräldralös flicka vid handen var på väg tillbaka till sitt
kloster.

%%K t14-tit-10 sub
\VUpk{10.2pt}{40}{Under vår hemfärd.}
\VUsaut{3.0mm}

%%K t14-20-1 p
20. --- I det ögonblick då min far åter slöt sig till oss, skrattade Johannes högt, när han
såg det smutsiga ansiktet och den underliga trasiga klädseln hos två \VUgras{sotare}
\textsuperscript{(91)} som var alldeles sotiga.

%%K t14-21-1 p
21. --- Jag ville köpa en växt åt min mor hos \VUgras{blomsterförsäljerskan}
\textsuperscript{(92)}, men min far gjorde mig uppmärksam på, att den skulle vara mycket tung
att bära med sig och att det vore bättre att skaffa
\VUgras{apelsiner}. Vi närmade oss \VUgras{försäljerskan}
\textsuperscript{(94)} och när \VUgras{guvernanten} \textsuperscript{(93)}, som höll på att
välja åt sin skyddsling, hade avslutat sitt köp, lät vi betjäna oss.

%%K t14-22-1 p
22. --- När vi vände hem igen, mötte vi flera bekanta : en gammal väninna till min familj, som
stödde armen på sin \VUgras{sällskapsdam} \textsuperscript{(95)}, \VUgras{ingenjören}
\textsuperscript{(96)} för broar och vägar, och en av mina kusiner med sin
\VUgras{barnsköterska} \textsuperscript{(98)}.

%%K t14-22-2 p
Sedan mötte vi min mors \VUgras{modist} \textsuperscript{(97)} och två
\VUgras{munkar} \textsuperscript{(99)} som var främmande för staden, vilka kom
fram till oss för att fråga min far om vägen till stationen.

\VUsaut{6.0mm}
\VUfilet{22mm}
